%%%%%%%%%%%%%%%%%%%%%%%%%%%%%%%%%%%%%%%%%
% Medium Length Professional CV
% LaTeX Template
% Version 2.0 (8/5/13)
%
% This template has been downloaded from:
% http://www.LaTeXTemplates.com
%
% Original author:
% Trey Hunner (http://www.treyhunner.com/)
%
% Important note:
% This template requires the resume.cls file to be in the same directory as the
% .tex file. The resume.cls file provides the resume style used for structuring the
% document.
%
%%%%%%%%%%%%%%%%%%%%%%%%%%%%%%%%%%%%%%%%%

%----------------------------------------------------------------------------------------
%	PACKAGES AND OTHER DOCUMENT CONFIGURATIONS
%----------------------------------------------------------------------------------------

\documentclass{resume} 
\usepackage[left=0.75in,top=0.6in,right=0.75in,bottom=0.6in]{geometry}
\usepackage{hyperref}
\usepackage [spanish]{babel}
\usepackage{picins}
\usepackage{graphicx}
\usepackage[utf8]{inputenc}
\usepackage{xcolor}
\newcommand{\tab}[1]{\hspace{.2667\textwidth}\rlap{#1}}
\newcommand{\itab}[1]{\hspace{0em}\rlap{#1}}
\name{Mar\'{i}a Jos\'{e} Bustamante Rosell}
\address{(+1) 512-576-3501 \\ \href{mailto:mbustamante@fisk.edu}{mbustamante@fisk.edu}}
\newcommand{\apj}{ApJ}
\newcommand{\apjl}{ApJL}
\newcommand{\prl}{PRL}
\newcommand{\mnras}{MNRAS}

\begin{document}

%----------------------------------------------------------------------------------------
%	EDUCATION SECTION
%----------------------------------------------------------------------------------------

%\begin{rSection}{}

%\end{rSection}

\begin{rSection}{Positions Held}\vspace{0.4cm}
{\bf Scientist} \hfill {\em March 2026-Present}
\\ Universities Space Research Association (USRA)

{\bf EMIT Postdoctoral Fellow} \hfill {\em 2024-2026}
\\ Fisk and Vanderbilt Universities, Nashville

{\bf Postdoctoral Scholar} \hfill {\em 2021-2024}
\\ University of California, Santa Cruz
\end{rSection}
\begin{rSection}{Education}\vspace{0.4cm}


{\bf The University of Texas at Austin} \hfill {\em 2014-2021} 
\\ PhD in Physics, Advisor: Aaron Zimmerman
\\

{\bf Pontificia Universidad Cat\'{o}lica del Per\'{u} (PUCP)} \hfill {\em 2008-2013} 
\\ BSc in Physics (first in class), 2012
\end{rSection}

%----------------------------------------------------------------------------------------
%   RESEARCH FUNDING SUMMARY
%----------------------------------------------------------------------------------------
\begin{rSection}{Research Funding Summary}\vspace{0.4cm}
Total competitive funding secured as PI: \$588,370
\end{rSection}

%----------------------------------------------------------------------------------------
%   SPACE & OBSERVATORY PROGRAMS
%----------------------------------------------------------------------------------------
\begin{rSection}{Space \& Observatory Programs}\vspace{0.4cm}
{\bf PI — JWST Cycle 3 Archival Research Program ID 5355} \hfill {\em 2025--2028}\\
Space Telescope Science Institute / JWST (NASA/ESA). Award: \$220,982. \href{https://www.stsci.edu/jwst/science-execution/program-information?id=5355}{Program page}.\\
{\it Title:} Adapting the Constrained Matrix Factorization Deblender \texttt{Scarlet} to Separate Overlapping Sources in JWST IFU Data.

{\bf Co-I — JWST General Observer Program ID 3805 (Cycle 2)} \hfill {\em 2023--2025}\\
Space Telescope Science Institute / JWST (NASA/ESA). \href{https://www.stsci.edu/jwst/science-execution/program-information?id=3805}{Program page}.\\
{\it Contribution:} Proposal design, observing plan, analysis strategy.
\end{rSection}

%----------------------------------------------------------------------------------------
%   NASA LISA PREPARATORY SCIENCE (LPS)
%----------------------------------------------------------------------------------------
\begin{rSection}{NASA LISA Preparatory Science (LPS)}\vspace{0.4cm}
{\bf PI — Electromagnetically-Driven GW Searches for EMRIs} \hfill {\em 2025--2028}\\
Proposal 24-LPS24-0023. Award: \$367,388.
Multi-messenger Bayesian framework linking EM timing (QPE/QPO) with GW morphology for ``wet'' EMRIs.

{\bf Co-I — A Robust Search for Faint and Weakly Modeled Sources in LISA} \hfill {\em 2024}\\
Semi-coherent, weakly modeled pipeline; chirplet banks + HMM/PCA; {\it includes GPU-accelerated template evaluation}. NASA panel rating: {\bf Very Good}.

{\bf Science PI — LISA Sources in the Galactic Center (Proposal 22-LPS22-0035)} \hfill {\em 2022}\\
Tidally-aware X-MRI waveform development; GLASS integration; NASA panel rating: {\bf Excellent / Very Good}.
\end{rSection}

\newpage

\begin{rSection}{Collaborations}\vspace{0.4cm}         
 LISA Consortium \hfill {\em 2021-Present}\\
 Young Supernova Experiment (YSE) \hfill {\em 2021-2023}\\
 LIGO-HET Response (LIGHETR)\\ 
 Electromagnetic Counterpart Search \textbf{Team Leader} \hfill {\em 2019-2021}\\
 SXS Collaboration \hfill {\em 2018-2021}\\
 LIGO/Virgo Scientific Collaboration \hfill {\em 2018-2021}\\
 MEDIDO PROJECT \hfill {\em 2016-2018}\\
 MINERvA \hfill {\em 2013-2014}
\end{rSection}



\begin{rSection}{Fellowships and Awards}
\item Establishing Multimessenger Astronomy Interdisciplinary Training (EMIT) \\Postdoctoral Fellowship \textit{Full financial support} \hfill{\em2024-2027}
\item Good Neighbour Scholarship \textit{Full year tuition} \hfill{\em2017-2018}
\item John A. Wheeler Fellowship in Physics \textit{Summer tuition and stipend}\hfill{\em Summer 2016}
\item Full Scholarship for pursuing the MSc in Physics at PUCP. \hfill{\em2012 }
\end{rSection}

\begin{rSection}{Mentorship and Leadership}


{\bf Founder \& Director — Fisk–Vanderbilt IDEA Summer Scholars Program} \hfill {\em 2023--Present}\\
Created a fully funded summer program pairing three Fisk data-science undergraduates with Vanderbilt postdocs on astrophysics/data-science projects (selection, matching, funding).
\end{rSection}

\begin{rSection}{Teaching and Mentoring}
{\bf Co-Director, IDEA Summer Scholars Program (Fisk–Vanderbilt)} \hfill {\em 2024–Present}\\
Designed and co-led a bridge research and training program for undergraduates from HBCUs entering astrophysics. Supervised 6 mentees; coordinated joint mentoring with Vanderbilt and UT Austin faculty.\\[4pt]
{\bf Instructor / Teaching Assistant}, University of Texas at Austin \hfill {\em 2014–2019}\\
Led discussion and lab sections for Classical Mechanics, Introductory Astronomy, and Computational Physics. Developed assessment rubrics and mentoring strategies for diverse student backgrounds.\\[4pt]
{\bf Graduate Mentor}, Vanderbilt {\em 2025–2024}\\
Mentored one graduate student in gravitational-wave data analysis and high-performance computing.\\
{\bf Undergraduate Mentor}, \\
\end{rSection}

\begin{rSection}{Teaching Experience}

\begin{rSubsection}{UT-Austin}{\em 2014-2021}{}{}
\item TA for General Relativity
\item AI for Physical Science 303  (Instructor of record for the class)
\item TA for Popular Astronomy
\item TA for General Physics I
\end{rSubsection}
\begin{rSubsection}{PUCP}{\em 2011-2012}{}{}
\item TA for Cosmology
\item TA for Introduction to Research
\end{rSubsection}
%------------------------------------------------

\end{rSection}
%	EXAMPLE SECTION
%----------------------------------------------------------------------------------------

\begin{rSection}{Selected Talks and Colloquia}
\item  Duke University Cosmology Seminar
\subitem{Science with the millihertz Gravitational Wave Sky: Gravitational Wave Timing Arrays and Electromagnetically Driven EMRI Searches} \textit{Durham}\hfill{\em April 4 2025}
\item University of Washington Astronomy Colloquium
\textit{Seattle}\hfill{\em February 15 2024}
\subitem{Black Holes and How to Find Them}
\item Space Science Seminar
\textit{NASA Marshall}\hfill{\em March 7 2023}
\subitem{Kinematic Constraints on Supermassive Black Holes in the Local Group and Beyond}
\item Rochester Institute of Technology Astronomy Colloquium
\textit{Rochester}\hfill{\em December 5 2022}
\subitem{The supermassive black hole at the center of Leo I}
\item Lisa Symposium XIV \textit{Virtual} \hfill{\em July 25 2022}
\subitem{Gravitational Wave Timing Array}
\item Intermediate-Mass Black Holes Conference \textit{Puerto Rico}\hfill{\em March 1 2022}
\subitem{Invited Talk - Dynamical analysis of the matter content in the Dwarf Spheroidal Galaxy Leo I }
\item APS April Meeting 2020 \textit{Virtual}
\subitem The LIGO HET Response (LIGHETR) Project to Discover and Spectroscopically Follow Optical Transients Associated with Neutron Star Mergers
\item CENTRA seminar series \textit{Lisbon}\hfill{\em Dec 10 2020}
\subitem 
Self-torque and frame nutation in binary black hole simulations
\item DESY Theoretical BS Astroparticle Seminar, DESY Zeuthen, Germany \hfill{\em Mar 17 2020}
\item Brown Bag Seminar, The University of Texas at Austin, \textit{Austin}
\subitem Black Hole Spin Precession in the Self-Force Approximation
\item APS April Meeting 2019 \textit{Colorado}
\subitem Dynamical analysis of the dark matter and black hole mass in the dwarf spheroidal LEO I
\item EXGal Seminar, The University of Texas at Austin, \textit{Austin}
\subitem Dark matter density profiles in local dSphs: \hfill{\em November 30, 2017}
\subitem The Leo I case
\item OPINAS Seminar, Max-Planck-Institut f\"ur extraterrestrische Physik, \textit{Garching}
\subitem LEO I'S CENTRAL POTENTIAL: \hfill{\em July 31, 2017}
\subitem Measurement difficulties and estimations from dynamical models
\item OPINAS Ringberg Meeting, Max-Planck-Institut f\"ur extraterrestrische Physik, \textit{Schloss Ringberg}
\subitem Central Dark Matter Density Profiles \hfill{\em April 6, 2017}
\subitem LEO I
\end{rSection}


\begin{rSection}{Publications} 
\noindent
\begingroup
\renewcommand{\section}[2]{}%
\nocite{*}
\bibliography{refs}{}
\bibliographystyle{plainyr-rev}
\endgroup
\end{rSection}
%\newpage

\begin{rSection}{Invitations and Visits}
\item \textit{Visiting Scientist at Max-Planck-Institut f\"ur extraterrestrische Physik (Medido Project)} \hfill{\em 2017}
\subitem - Collaborated on the kinematic measurements and dynamical models for \subitem Dwarf Spheroidals as part of the Medido Project.
\subitem - Planning of observations of ultra faint Dwarf Spheroidals with VIRUS-W.
\item \textit{Visiting Scientist at Fermilab (MINERvA)}  \hfill{\em 2014-2015}
\subitem - Leaded the implemention of the fourth Muon Monitor in the NuMI beamline.
\subitem - Studied and calibrated the flux pipeline in the NuMI beamline.
\subitem - Work on the study and design of new Cherenkov detectors for DUNE.
\end{rSection}

\begin{rSection}{Professional Service}
\item \textbf{Scientific Organizing Committee}: LISA Symposium \hfill{\em 2026}
\item \textbf{Proposal Reviewer}: NASA Euclid General Investigator Program (EGIP) \hfill{\em 2025}
\item \textbf{Local Organizer}: 3rd LISA Sprint, Huntsville, AL \hfill{\em 2025}
\item \textbf{Peer Reviewer}: Nature Astronomy \hfill{\em 2023}


\end{rSection}
%\begin{rSection}{Observation Time Granted}

%\item Granted a total of 10 days of time on McDonald Observatory's 2.7m to observe Leo I as part of the Medido Project
%\subitem Used the IFU spectrograph VIRUS-W and reduced the data with Panacea
%\item Granted 60 hr of priority-zero time for kilonovae follow-up on the Hobby-Eberly Telescope.
%\subitem Used the IFU spectrographs VIRUS and LRS2. Worked on kilonova detection pipeline (see section of research statement on SCARLET).
%\end{rSection}
\newpage
\begin{rSection}{Activities} \itemsep -3pt
\item Research Assistant for LIGO and SXS collaborations \hfill{\em 2018-2021}
\item Research Assistant for HETDEX collaboration \hfill{\em 2015-2021}
\item Research Assistant in PUCP's High Energy Physics Group \hfill{\em 2013-2014}
\item Research Assistant in PUCP's Quantum Optics Group \hfill{\em 2010-2012}
\item Volunteer Teacher for NGO ``Bruce Peru" \hfill{\em 2010-2012}
\item Research Assistant for NGO Otra Mirada \hfill{\em 2009}
\item Development and Implementation of Online Queries -- Overseas Development Institute (ODI) \hfill{\em 2008}
\item Member of PUCP's Complex Systems Investigation Group \hfill{\em 2007-2012}
\item Foundational Member of Astronomy Group in PUCP \hfill{\em 2006}
\item Volunteer Firefighter \hfill{\em 2004-2006}

\end{rSection}

\begin{rSection}{Outreach} \itemsep -2pt
\item Fisk-Vanderbilt Interdisciplinary Data-science Experience in Astrophysics (IDEA) Summer Scholarship --
\textit{Nashville 2025.}
Program organizer and mentor
\item Society for Advancement of Chicanos/Hispanics \& Native Americans in Science (SACNAS) --
\textit{Portland 2023.}
Panelist in a session on Multi-messenger Astronomy.
\item Lamat Program --
\textit{Santa Cruz 2023.}
Mentor for Undergraduate Summer Student.
\item EMIT Summer School in Multimessenger Astronomy -- \textit{Vanderbilt 2023.}
Developed tutorial notebooks for parameter estimation of gravitational waves.
\item Conferences for Undergraduate Women in Physics (CUWiP) -- 
\textit{Santa Cruz 2023.} Volunteer
\item Gravitational Wave Astronomy Northwest meeting (GWANW) --
\textit{Ligo Hanford 2022.} Lisa tutorial
\item Schrodinger's Pint -- \textit{Austin 2022.} Public lecture ``Black holes and where to find them"
\item Astronomy on Tap talk -- \textit{Austin 2020.}
News section
\item Nerd Nite -- \textit{Austin 2020.}
Public lecture about gravitational waves.
\item Black Hole Night -- \textit{Austin 2019.}
Hosted event to promote collaboration between local artists and scientists on topics relating to black holes.
\item CASA BAGRE - \textit{Lima 2019.}
Talk on Colliding Black Holes (``Agujeros Negros Collisionando") in collaboration with audiovisual artists.
\item Astronomy on Tap talk -- \textit{Austin 2017.} 
Public lecture about cosmic inflation.
\item Girl Day -- \textit{2019}. Volunteer for a UT event focused on exposing girls to STEM.
\item Alice in Wonderland volunteer -- program focused on introducing high school girls to physics research.
\end{rSection}

\end{document}
